\SourceLeaf{226}%
\noindent Ceterum quum inde a puncto \(A\) bini rami utriusque circuli maximi proficiscantur, duo quidem ibi anguli formantur, quorum alter alterius complementum ad \(180^0\): sed analysis nostra monstrat, eos ramos adoptandos esse, quorum directiones cum sensu progressionis a puncto \(L\) ad \(L'\), et a puncto \(L''\) ad \(L'''\) consentiunt: quibus intellectis simul patet, quum circuli maximi duobus punctis concurrant, arbitrarium esse, utrum eligatur. Loco anguli \(A\) etiam arcus inter polos circulorum maximorum, quorum partes sunt arcus \(LL',L''L'''\), adhiberi potest: manifesto autem polos tales accipere oportet, qui respectu horum arcuum similiter iacent, puta vel uterque polus ad dextram iacens, dum a \(L\) versus \(L'\) atque ab \(L''\) versus \(L'''\) procedimus, vel uterque ad laevam.

VII. Sint \(L,L',L''\) tria puncta in superficie sphaerica, statuamusque brevitatis caussa
\[
\begin{aligned}
\cos(1)L&=x,&\cos(2)L&=y,&\cos(3)L&=z\\
\cos(1)L'&=x',&\cos(2)L'&=y',&\cos(3)L'&=z'\\
\cos(1)L''&=x'',&\cos(2)L''&=y'',&\cos(3)L''&=z''
\end{aligned}
\]
\noindent nec non
\[
xy'z''+x'y''z+x''yz'-xy''z'-x'yz''-x''y'z=\Delta
\]
\noindent Designet \(\lambda\) polum circuli maximi, cuius pars est arcus \(LL'\), et quidem eum, qui respectu huius arcus similiter iacet, ac punctum \((1)\) respectu arcus \((2)(3)\). Tunc erit, ex theoremate praecedente, \(yz'-y'z=\cos(1)\lambda\mathord{.}\sin(2)(3)\mathord{.}\sin LL'\), sive, propter \((2)(3)=90^0\),
\[
\begin{aligned}
yz'-y'z&=\cos(1)\lambda\mathord{.}\sin LL',\ \text{et perinde}\\
zx'-z'x&=\cos(2)\lambda\mathord{.}\sin LL'\\
xy'-x'y&=\cos(3)\lambda\mathord{.}\sin LL'
\end{aligned}
\]
\noindent Multiplicando has aequationes resp. per \(x'',y'',z''\) et addendo, obtinemus adiumento theorematis secundi in \(V\) prolati
\[
\Delta=\cos\lambda L''\mathord{.}\sin LL'
\]
\noindent Iam tres casus sunt distinguendi. \textit{Primo}, quoties \(L''\) iacet in eodem circulo maximo, cuius pars est arcus \(LL'\), erit \(\lambda L''=90^0\), adeoque \(\Delta=0\). Quoties vero \(L''\) iacet extra circulum illum maximum, aderit casus \textit{secundus}, si est ab eadem parte, a qua est \(\lambda\), \textit{tertius}, si ab opposita: in his casibus puncta \(L,L',L''\)
